% figures/index-layout.tex -- sequential index layout and posting structure.
% Pure TikZ diagram: field names and byte widths come from include/index.hpp
% and the v1 serializer in src/index.cpp (28 B header, 12+len doc entry,
% 28+len term entry, 8+4*tf posting entry). No measured values.
\begin{figure}[!htbp]
\centering
\begin{tikzpicture}[
  font=\footnotesize,
  >=Stealth,
  x=1cm, y=1cm,
  seglabel/.style={font=\footnotesize, text=mseblue, anchor=east},
  strip/.style={draw=black!50, rounded corners=1pt, line width=0.5pt,
                fill=black!4, anchor=west, align=left, inner sep=2.2pt,
                font=\scriptsize, text width=12.35cm},
  cell/.style={draw=black!55, rounded corners=1pt, line width=0.45pt,
               fill=white, align=center, inner sep=1.5pt,
               minimum height=0.52cm, font=\scriptsize},
  structbox/.style={draw=mseblue!70, rounded corners=3pt, line width=0.6pt,
                    fill=mseblue!4},
  title/.style={font=\small, text=black!80, anchor=west, align=left,
                text width=15.2cm},
  note/.style={font=\scriptsize, text=msegray, anchor=west, align=left},
]

% ======================================================= (a) sequential layout
\node[title] at (0.20,7.55)
  {(a) Sequential layout of \texttt{index.bin} (file order; lengths not to
   scale)};
\node[note, text width=15.2cm] at (0.20,7.10)
  {v2 (\texttt{-{}-compress}) changes only the posting segment: document IDs
   and positions become delta-coded varints, \texttt{version} is set to 2, and
   the other segments are unchanged.};

% --- header
\node[seglabel] at (2.42,6.30) {header (28 B)};
\node[strip] at (2.55,6.30)
  {Header, 28 B: magic \texttt{"MSE1"}, then five count and flag fields\\
   \texttt{version} u32 · \texttt{flags} u32 · \texttt{docCount} u32 ·\\
   \texttt{termCount} u32 · \texttt{tokenCount} u64};
% --- document table
\node[seglabel] at (2.42,5.35) {doc table};
\node[strip] at (2.55,5.35)
  {12 + len(title) B each, in increasing docID order\\
   \texttt{docID} u32 · \texttt{keptTokens} u32 · \texttt{titleLen} u32 ·
   \texttt{title} (variable bytes)};
% --- term table
\node[seglabel] at (2.42,4.40) {term table};
\node[strip] at (2.55,4.40)
  {28 + len(term) B each, in lexicographic term order
   (\texttt{offset} points to the posting block)\\
   \texttt{len} u32 · \texttt{term} · \texttt{df} u32 · \texttt{tfTotal} u64 ·
   \texttt{offset} u64 · \texttt{bytes} u32};
% --- postings
\node[seglabel] at (2.42,3.45) {postings};
\node[strip] at (2.55,3.45)
  {One block per term, 8·$\Sigma$df + 4·$\Sigma$tf B, no padding between
   blocks\\
   \texttt{df} × (\texttt{docID} u32 · \texttt{tf} u32) + $\Sigma$tf
   \texttt{positions} u32};

% ======================================================= (b) term -> block
\node[title] at (0.20,2.95)
  {(b) One term record $\rightarrow$ its posting block};

% --- term record: six field cells, byte width on the second line
\node[cell, minimum width=0.90cm] (len) at (0.85,2.20) {len\\[-1pt]{\scriptsize u32}};
\node[cell, minimum width=1.60cm] (term) at (2.15,2.20) {term\\[-1pt]{\scriptsize var. bytes}};
\node[cell, minimum width=0.90cm] (df) at (3.50,2.20) {df\\[-1pt]{\scriptsize u32}};
\node[cell, minimum width=1.30cm] (tftotal) at (4.70,2.20) {tfTotal\\[-1pt]{\scriptsize u64}};
\node[cell, minimum width=1.20cm] (offset) at (6.00,2.20) {offset\\[-1pt]{\scriptsize u64}};
\node[cell, minimum width=1.10cm] (bytes) at (7.20,2.20) {bytes\\[-1pt]{\scriptsize u32}};

\node[note, text width=7.90cm] at (0.35,1.48)
  {One term record = 28 + len(term) B; \texttt{offset} points to the first byte
   of that term's posting block.};

% --- posting block (two records shown; all fields are u32)
\draw[structbox] (9.55,0.38) rectangle (15.42,2.80);

\node[cell, minimum width=0.95cm] at (10.22,2.44) {docID};
\node[cell, minimum width=0.55cm] at (11.18,2.44) {tf};
\node[cell, minimum width=0.72cm] at (12.00,2.44) {pos$_0$};
\node[cell, minimum width=0.72cm] at (12.85,2.44) {pos$_1$};
\node[cell, minimum width=0.55cm] at (13.60,2.44) {$\cdots$};
\node[cell, minimum width=1.35cm] at (14.68,2.44) {pos$_{\mathrm{tf}-1}$};
\node[font=\scriptsize, text=msegray] at (12.54,2.02) {$\vdots$};
\node[cell, minimum width=0.95cm] at (10.22,1.60) {docID};
\node[cell, minimum width=0.55cm] at (11.18,1.60) {tf};
\node[cell, minimum width=0.72cm] at (12.00,1.60) {pos$_0$};
\node[cell, minimum width=1.15cm] at (13.30,1.60) {$\cdots$};
\node[cell, minimum width=1.35cm] at (14.68,1.60) {pos$_{\mathrm{tf}-1}$};

\node[note, anchor=north west, text width=5.40cm] at (9.76,1.24)
  {Each entry is 8 + 4·tf B: \texttt{docID}, \texttt{tf},\\
   then tf ascending u32 positions; df entries, packed.};

% --- term record -> posting block
\draw[->, line width=0.6pt, draw=black!65]
  (bytes.east) -- (9.55,2.20);
\node[note, anchor=south, inner sep=1pt] at (8.55,2.26) {\texttt{seek(offset)}};

\end{tikzpicture}
\caption{Index layout: \texttt{index.bin} stores the header, document table,
term table, and posting segment in order. A term record points to its posting
block with \texttt{offset} and \texttt{bytes}. Within a block, each posting
stores the document ID and term frequency followed by tf ascending positions,
so a phrase query only needs to compare adjacent positions.}
\label{fig:index-layout}
\end{figure}
